The third example demonstrates the pipeline's relevance gate: when no
candidate ad crosses the minimum similarity threshold, the system returns a
purely helpful response with no sponsored content.
Figure~\ref{fig:example-3} shows the result for a health-and-symptoms query
--- an inherently non-commercial topic for which no ad in the inventory is
relevant.

\begin{figure}[H]
    \centering
    \includegraphics[width=\textwidth]{images/illustrative_example_2.png}
    \caption{Chat response to a health-symptoms query.}
    \label{fig:example-3}
\end{figure}

\subsubsection{Stage 1: Context Extraction}
\label{subsubsec:example3-context}

The context-extraction module classifies the query as a general information
request with no commercial intent:

\begin{lstlisting}[style=jsonstyle, caption={Context extraction output for the
health query.}, label={lst:example3-context}]
{
  "intent":          "general_question",
  "ad_receptivity":  "low",
  "purchase_signal": 0.00,
  "topics":          ["health", "symptoms", "medical"],
  "entities":        []
}
\end{lstlisting}

\noindent
A \texttt{purchase\_signal} of~0.00 and \texttt{ad\_receptivity} of
\texttt{low} signal that the user is not in a buying mindset.  Unlike
\texttt{none} (which would skip retrieval entirely), \texttt{low} still
permits a vector search --- but the re-ranker applies no purchase-proximity
boost, and the similarity threshold is enforced strictly.

\subsubsection{Stage 2: Vector Search and Re-ranking}
\label{subsubsec:example3-retrieval}

The topics \emph{``health symptoms medical''} are encoded and searched against
the ad inventory.  Three candidates are returned, none of which are
thematically related to medical queries:

\begin{lstlisting}[style=jsonstyle, caption={All three vector-search candidates,
all below the 0.35 similarity threshold.}, label={lst:example3-retrieval}]
[
  { "product_name": "Dyson V15 Detect",     "similarity": 0.2943 },
  { "product_name": "Nike Pegasus 41",      "similarity": 0.2307 },
  { "product_name": "Allbirds Tree Runners","similarity": 0.1782 }
]
\end{lstlisting}

\noindent
Every score falls below the minimum threshold of~0.35.  The re-ranker
therefore discards all three candidates and forwards an empty ad list to
the synthesis stage.

\subsubsection{Stage 3: Response Generation (No Sponsored Content)}
\label{subsubsec:example3-response}

With an empty ad list, the orchestrator assembles a system prompt that
contains no \texttt{<sponsored\_context>} block.  The LLM produces a
straightforward medical-information response --- identifying potential causes
(viral or bacterial infection, impetigo, hand-foot-and-mouth disease) and
recommending a doctor consultation --- without any product mention or
disclosure label.  No impression event is written and no campaign spend is
recorded for this turn.
